% ****************************************************************************************************
% classicthesis-config.tex
% Adapted from the IMISE classicthesis template (https://github.com/IMISE/imise-classicthesis),
% itself based on classicthesis by Andre Miede and Ivo Pletikosic.
% Adapted for the book "The Cosmic Mass" by Mostafa Derraz.
% ****************************************************************************************************

% ****************************************************************************************************
% 0. Encoding (and color names used by classicthesis, before any package loads xcolor)
% ****************************************************************************************************
\PassOptionsToPackage{dvipsnames,usenames}{xcolor}
\PassOptionsToPackage{utf8}{inputenc}
  \usepackage{inputenc}
\PassOptionsToPackage{T1}{fontenc}
  \usepackage{fontenc}

% ****************************************************************************************************
% 1. classicthesis options
% ****************************************************************************************************
\PassOptionsToPackage{
  drafting=true,             % prints the version line at the bottom of pages (set to false for the final book)
  tocaligned=false,
  dottedtoc=false,
  eulerchapternumbers=false,
  floatperchapter=true,      % Figure 3.1, Table 3.1, ...
  eulermath=false,
  beramono=false,            % beramono is not installed here
  palatino=true,
  style=classicthesis
}{classicthesis}

% ****************************************************************************************************
% 2. Book data (edit here)
% ****************************************************************************************************
\newcommand{\myTitle}{The Cosmic Mass\xspace}
\newcommand{\mySubtitle}{One Number for Mass and Charge\xspace}
\newcommand{\myVolume}{Volume 1\xspace}
\newcommand{\myName}{Mostafa Derraz\xspace}
\newcommand{\myAffiliation}{Independent Researcher\xspace}
\newcommand{\myORCID}{0000-0002-2307-9298\xspace}
\newcommand{\myEmail}{m.derraz@proton.me\xspace}
\newcommand{\myLocation}{Morocco\xspace}
\newcommand{\myTime}{2026\xspace}
\newcommand{\myVersion}{First draft\xspace}

% ****************************************************************************************************
% 3. Packages
% ****************************************************************************************************
\PassOptionsToPackage{american}{babel}
\usepackage{babel}
\usepackage{csquotes}

\usepackage[round,authoryear]{natbib}

\PassOptionsToPackage{fleqn}{amsmath}
  \usepackage{amsmath}
\usepackage{amssymb}
\usepackage{amsthm}
\usepackage{mathtools}

\usepackage{graphicx}
\usepackage{scrhack}
\usepackage{xspace}
\PassOptionsToPackage{activate={true,nocompatibility},final,tracking=true,kerning=true,spacing=true,factor=1100,stretch=10,shrink=10}{microtype}
\usepackage{microtype}
\usepackage{siunitx}
\usepackage{booktabs}
\usepackage{tabularx}
\usepackage{longtable}   % the symbol list of Appendix D, which runs over two pages
\setlength{\extrarowheight}{3pt}
\newcommand{\tableheadline}[1]{\multicolumn{1}{l}{\spacedlowsmallcaps{#1}}}
\newcommand{\myfloatalign}{\centering}
\setlength{\emergencystretch}{2.5em}
\usepackage{enumitem}
\setlist{itemsep=0.25em}

\usepackage{tikz}
\usetikzlibrary{arrows.meta,calc,angles,quotes,decorations.pathreplacing,patterns}
\usepackage{pgfplots}
\pgfplotsset{compat=1.16}
\usepackage[most]{tcolorbox}

% ****************************************************************************************************
% 4. Her Majesty herself
% ****************************************************************************************************
\PassOptionsToPackage{hyperfootnotes=false}{hyperref}
\usepackage{classicthesis}

% Version line printed at the foot of each page while drafting (no dash)
\ifdefined\PrelimText\renewcommand{\PrelimText}{\footnotesize[\,\myVersion, \today\,]}\fi

\hypersetup{%
  colorlinks=true, linktocpage=true, pdfstartpage=1, pdfstartview=FitV,%
  breaklinks=true, pageanchor=true,%
  pdfpagemode=UseOutlines,%
  plainpages=false, bookmarksnumbered, bookmarksopen=true, bookmarksopenlevel=1,%
  hypertexnames=true, pdfhighlight=/O,%
  urlcolor=CTurl, linkcolor=CTlink, citecolor=CTcitation,%
  pdftitle={\myTitle, \myVolume: \mySubtitle},%
  pdfauthor={\myName},%
  pdfsubject={General Sinus, General Numbers, Cosmic Mass},%
  pdfkeywords={General Numbers, Composer Unit, Cosmic Mass, General Sinus, mass, electric charge},%
  pdfcreator={pdfLaTeX},%
  pdfproducer={LaTeX with hyperref and classicthesis}%
}

\usepackage{cleveref}

% ****************************************************************************************************
% 5. Theorem-like statements
%    Axioms have their own counter; theorems, definitions and remarks share one counter,
%    both numbered within the chapter. Within Chapters 4 and 6 this keeps the order and numbering of the
%    original papers, with the chapter number raised by one since the new Chapter 2:
%    Axiom 3.1, Theorem 5.6, Remark 5.15 of the papers are Axiom 4.1, Theorem 6.6, Remark 6.15 here.
% ****************************************************************************************************
\theoremstyle{plain}
\newtheorem{theorem}{Theorem}[chapter]
\newtheorem{axiom}{Axiom}[chapter]
\newtheorem{corollary}[theorem]{Corollary}
\newtheorem{proposition}[theorem]{Proposition}
\theoremstyle{definition}
\newtheorem{definition}[theorem]{Definition}
\theoremstyle{remark}
\newtheorem{remark}[theorem]{Remark}

% ****************************************************************************************************
% 6. Notation used throughout the book
% ****************************************************************************************************
\newcommand{\pb}{\bar{\phi}}                        % projection angle of the composer axis
\newcommand{\Gp}{\mathcal{G}_{\bar{\phi}}}          % Composer Unit
\newcommand{\Ap}{\mathcal{A}_{\bar{\phi}}}          % amplitude of the Composer Unit
\newcommand{\up}{\vec{u}_{\bar{\phi}}}              % direction of the Composer Unit
\newcommand{\GN}{\mathbb{G}}                        % General Number Plane
\newcommand{\GNp}{\mathbb{G}_{\bar{\phi}}}          % Relative Plane
\newcommand{\xv}{\vec{x}}
\newcommand{\yv}{\vec{y}}
\newcommand{\er}{\vec{e}_r}
\newcommand{\R}{\mathbb{R}}
\newcommand{\C}{\mathbb{C}}
\newcommand{\Z}{\mathbb{Z}}
\newcommand{\seg}[1]{[#1]}                          % length of a segment, as in the papers
\newcommand{\Zb}{\bar{Z}}                            % projected image of the point Z

% ****************************************************************************************************
% 7. Boxes and small layout helpers
% ****************************************************************************************************
\definecolor{plainbg}{RGB}{241,245,250}
\definecolor{plainframe}{RGB}{95,120,160}
\definecolor{notebg}{RGB}{250,246,238}
\definecolor{noteframe}{RGB}{170,120,60}

% "In plain words": an explanation in everyday language (never numbered)
\newtcolorbox{plainwords}[1][In plain words]{%
  enhanced, breakable, colback=plainbg, colframe=plainframe, boxrule=0pt, leftrule=2.5pt,
  arc=0pt, outer arc=0pt, left=8pt, right=8pt, top=6pt, bottom=6pt,
  fonttitle=\normalfont\scshape, coltitle=plainframe, colbacktitle=plainbg,
  title={#1}, attach title to upper={\par\smallskip}, before skip=10pt, after skip=10pt,
  lines before break=4}

% "Worked example": a short calculation that uses the formulas exactly as stated (never numbered)
\newtcolorbox{workedexample}[1][Worked example]{%
  enhanced, breakable, colback=notebg, colframe=noteframe, boxrule=0pt, leftrule=2.5pt,
  arc=0pt, outer arc=0pt, left=8pt, right=8pt, top=6pt, bottom=6pt,
  fonttitle=\normalfont\scshape, coltitle=noteframe, colbacktitle=notebg,
  title={#1}, attach title to upper={\par\smallskip}, before skip=10pt, after skip=10pt,
  lines before break=4}

% Start a new page here if fewer than #1 lines are left on the current page; do nothing otherwise
\makeatletter
\newcommand{\keepblock}[1]{\par\begingroup\dimen@=\pagegoal\advance\dimen@-\pagetotal
  \ifdim\pagetotal>\z@\ifdim\dimen@<#1\baselineskip\newpage\fi\fi\endgroup}
\makeatother

% Place and date line that opens a historical scene (never left alone at the foot of a page)
\newcommand{\scene}[1]{\par\medskip\noindent{\centering\spacedlowsmallcaps{#1}\par}\nobreak\medskip}

% Author placeholders (easy to find and replace)
\newcommand{\placeholder}[1]{\textcolor{Maroon}{[#1]}}

% Chapter opening line
\newcommand{\chapterlead}[1]{\begingroup\itshape #1\par\endgroup\bigskip}

% Epigraph
\newcommand{\bookepigraph}[2]{%
  \begin{flushright}\begin{minipage}{0.78\textwidth}\raggedleft\itshape #1\\[0.4em]\normalfont\small #2\end{minipage}\end{flushright}\bigskip}
